\documentclass[../../main.tex]{subfiles}
\begin{document}

\begin{center}
  {\fontsize{48pt}{58pt}\selectfont \textbf{2006}\normalsize\textbf{年度}}
\end{center}

\begin{table}[htb]
  \centering
  \caption{出題テーマと難易度}
  \begin{tabular}{|c|p{9cm}|c|}
    \hline
    \textbf{問題番号} & \textbf{テーマ} & \textbf{難易度} \\
    \hline
    第1問 & 点の漸化式と双曲線・円周 & \\
    \hline
    第2問 & 記号の表示操作と確率 & \\
    \hline
    第3問 & 直線に関する対称移動と傾き & \\
    \hline
    第4問 & 整数解を持つ方程式と組の構成 & \\
    \hline
    第5問 & 漸化式で定めた数列と極限 & \\
    \hline
    第6問 & 指数関数の逆関数と定積分 & \\
    \hline
  \end{tabular}
\end{table}

\begin{table}[htb]
  \centering
  \caption{解答}
  \renewcommand{\arraystretch}{1.4}
  \begin{tabular}{|c|c|p{9cm}|}
    \hline
    \textbf{問題番号} & \textbf{小問} & \textbf{答え} \\
    \hline
    第1問 & (1) & $P_3$は曲線$xy=1$上にない（証明問題） \\
    \cline{2-3}
    & (2) & $P_4$も円周$x^2+y^2=1$上にある（証明問題） \\
    \hline
    第2問 & (1) & $P_2=(1-p)(2p^2-p+1)$ \\
    \cline{2-3}
    & (2) & $P_n=(1-p)p^{n-1}+(1-p)p^n
    +(n-1)(1-p)^3p^{n-2}$ \\
    \hline
    第3問 & (1) & $\tan\theta=
    \dfrac{(p-1)-(p+1)\alpha^2}{\alpha(\alpha^2+1-2p)}$ \\
    \cline{2-3}
    & (2) & $p=2$（存在の証明を含む） \\
    \hline
    第4問 & (1) & $(x,y,z)=(3,3,3),(3,3,6)$ \\
    \cline{2-3}
    & (2) & 条件を満たす$z$が存在する（証明問題） \\
    \cline{2-3}
    & (3) & 条件を満たす組は無数にある（証明問題） \\
    \hline
    第5問 & (1) & $b_n>2n$（証明問題） \\
    \cline{2-3}
    & (2) & $\displaystyle\lim_{n\to\infty}
    \dfrac{1}{n}\sum_{k=1}^{n}a_k=0$ \\
    \cline{2-3}
    & (3) & $\displaystyle\lim_{n\to\infty}na_n=\dfrac{1}{2}$ \\
    \hline
    第6問 & (1) & $f$は実数全体を定義域とする逆関数を持つ（証明問題） \\
    \cline{2-3}
    & (2) & $\displaystyle\int_8^{27}g(x)\,dx
    =39\log3-20\log2-12$ \\
    \hline
  \end{tabular}
\end{table}

\paragraph{試験形態}
6問．試験時間は要確認．

\end{document}
